\documentclass[10pt,twocolumn]{article}

\usepackage[margin=0.75in]{geometry}
\usepackage{amsmath}
\usepackage{booktabs}
\usepackage{graphicx}
\usepackage[hidelinks]{hyperref}
\usepackage{xcolor}

\newcommand{\blade}{\textsc{Blade}}
\newcommand{\wgpu}{\textsc{wgpu}}
\newcommand{\pending}[1]{\textcolor{red}{\textbf{RESULTS PENDING: #1}}}

\title{Do Fine-Grained GPU Resource States Pay for Themselves?\\
A Cross-Vendor Study with \blade}
\author{Author information to be finalized}
\date{Draft of July 2026}

\begin{document}
\maketitle

\begin{abstract}
Modern explicit graphics APIs make applications describe memory dependencies,
resource access scopes, and image layouts in detail. Portable APIs such as
WebGPU reconstruct and validate this state at run time, while render graphs
derive it from a declared workload. \blade{} explores a different point in the
design space: its Vulkan backend keeps ordinary images in
\texttt{GENERAL}, does not maintain per-resource states, and places a global
memory barrier at pass boundaries.

This paper asks when such coarse synchronization is competitive with precise
tracking and where it fails. We combine a controlled within-\blade{}
comparison of barrier placement with matched end-to-end \wgpu{} workloads,
CPU profiles, and captured Vulkan command streams. This design avoids
reimplementing a resource tracker in \blade{}, while making clear which
differences are causal and which belong to the complete abstraction. The
evaluation covers synthetic dependency graphs and representative compute and
graphics workloads on AMD and NVIDIA Vulkan implementations.
\pending{Replace this paragraph with quantitative findings, confidence
intervals, and the observed decision boundary before circulation.}
\end{abstract}

\noindent\colorbox{yellow!25}{%
\parbox{\dimexpr\columnwidth-2\fboxsep\relax}{%
\textbf{Working draft.} The methodology is specified, but result claims and
figures have deliberately not been filled with provisional measurements.}}

\section{Introduction}

Explicit GPU APIs expose synchronization so applications can preserve
parallelism and avoid unnecessary cache operations. That control has a cost:
the application or an abstraction layer must know how every buffer and image
subresource is used, maintain state across command buffers, validate conflicts,
and generate transitions. This machinery sits on a CPU-critical command
recording path. The \wgpu{} implementation describes its trackers as some of
the hottest code in its codebase~\cite{wgpu-trackers}.

\blade{}~\cite{blade} tests the opposite engineering hypothesis. Its public API has no
resource-state objects. On Vulkan, ordinary images remain in
\texttt{VK\_IMAGE\_LAYOUT\_GENERAL}; a broad memory barrier orders writes
before later reads and writes at pass boundaries. The design sacrifices
opportunities for overlap and resource-specific cache control in exchange for
a smaller API and negligible tracking work.

The question is newly relevant because
\texttt{VK\_KHR\_unified\_image\_layouts} guarantees, when enabled, that
\texttt{GENERAL} is as efficient as specialized layouts in most
uses~\cite{vulkan-unified-layouts}. The extension removes one reason for
tracking image state, but explicitly leaves barrier granularity as a separate
performance question.

This paper is organized around the following research questions:

\begin{description}
  \item[RQ1:] How does host cost scale with resources, passes, and dependency
  density for tracking-free and tracked command recording?
  \item[RQ2:] When do global barriers reduce device overlap or cause observable
  cache and pipeline stalls?
  \item[RQ3:] How do matched \blade{} and \wgpu{} implementations differ
  end-to-end, and which costs are visible in tracker profiles and generated
  Vulkan commands?
  \item[RQ4:] How do the answers differ between compute, graphics, and mixed
  workloads, and between AMD and NVIDIA architectures?
\end{description}

The intended contributions are:

\begin{itemize}
  \item a cost model that distinguishes tracking, barrier precision, barrier
  placement, and image layouts without claiming that every factor is
  independently controlled;
  \item a controlled within-\blade{} evaluation of automatic versus
  application-declared barrier placement;
  \item matched \blade{} and \wgpu{} workloads with profiles and Vulkan
  captures for end-to-end comparison;
  \item a reproducible cross-vendor evaluation over dependency-controlled
  microbenchmarks and application workloads; and
  \item practical guidance identifying the workload envelope in which
  tracking-free synchronization is appropriate.
\end{itemize}

\section{Background}

\subsection{Execution and memory dependencies}

A Vulkan barrier describes both an execution dependency between pipeline
stages and a memory dependency between access types. Resource-specific
barriers can additionally identify the affected buffer or image range.
Narrow scopes may preserve unrelated work, whereas overly broad scopes can
drain the pipeline~\cite{amd-barriers,nvidia-vulkan-tips}. A global memory
barrier is nevertheless attractive when many resources share the same
dependency because it combines them in one command and often maps to hardware
operations whose granularity is already coarse.

These dimensions must not be collapsed into a single ``barrier quality''
variable. A system may:

\begin{enumerate}
  \item place barriers only at true hazards or at every pass boundary;
  \item use global or resource-specific memory scopes;
  \item use broad or precise pipeline stages and access masks; and
  \item transition images between specialized layouts or retain one layout.
\end{enumerate}

\subsection{Image layouts}

Vulkan image layouts communicate usage information that may select compression
metadata, attachment paths, or transfer representations. Vendor guidance has
historically recommended specialized layouts, particularly on
AMD~\cite{amd-rdna-guide}. The unified-image-layouts extension reflects a
hardware and driver trend in which most layouts are physically identical. It
allows \texttt{GENERAL} in most positions and promises optimal performance,
while retaining exceptions such as initialization from
\texttt{UNDEFINED}, presentation, and some video interactions
~\cite{vulkan-unified-layouts}.

Consequently, this work uses ``no layout tracking'' to mean no steady-state
tracking for ordinary application images. It does not claim that all layout
transitions or image-specific synchronization disappear.

\subsection{WebGPU and \wgpu}

WebGPU defines usage scopes that reject conflicting uses and provide a safe,
portable programming contract~\cite{webgpu-spec}. \wgpu{} represents buffer
and texture states in dense trackers, merges per-scope usage, and generates
transitions between command buffers~\cite{wgpu-trackers}. This work does not
treat that work as accidental overhead: it implements guarantees that
\blade's unsafe API intentionally does not provide.

The comparison instead asks what those guarantees and their precision cost,
and whether a narrower native contract is useful for applications that already
know their dependency structure.

\section{\blade{} Synchronization Model}

\blade{} exposes transfer, compute, and render passes within a command encoder.
It does not attach a current state to resource handles. The Vulkan backend
places the following conceptual dependency before each pass:

\begin{verbatim}
srcStage  = ALL_COMMANDS
dstStage  = ALL_COMMANDS
srcAccess = MEMORY_WRITE
dstAccess = MEMORY_READ | MEMORY_WRITE
\end{verbatim}

Images are initialized from \texttt{UNDEFINED} to \texttt{GENERAL} and remain
in \texttt{GENERAL} for ordinary rendering, sampling, storage, and copies.
Presentation images still transition to the presentation layout. An encoder
option disables automatic pass barriers; the application may then call the
same global barrier only where its dependency graph requires one.

The evaluated contract is deliberately limited:

\begin{itemize}
  \item synchronization within one Vulkan queue and command stream;
  \item explicit host and cross-queue synchronization outside the pass model;
  \item explicit initialization and presentation transitions; and
  \item native validation during correctness testing, but no safe-API
  guarantee in production.
\end{itemize}

\section{Experimental Design}

\subsection{Configurations}

Table~\ref{tab:configs} separates the controlled \blade{} configurations from
the realistic tracked baseline. \texttt{B-explicit-all} checks that manual and
automatic insertion produce equivalent command streams. \texttt{W-wgpu} uses
the same shaders, resources, pass graph, and output checks, but remains an
end-to-end comparison because the APIs and runtime guarantees differ.

\begin{table*}[t]
\centering
\small
\begin{tabular}{lllll}
\toprule
ID & Tracking & Placement & Scope & Layout \\
\midrule
B-auto & none & every pass & global & general \\
B-hazard & none & hazards & global & general \\
B-explicit-all & none & every pass & global & general \\
W-wgpu & tracked & derived & resource & backend \\
\bottomrule
\end{tabular}
\caption{Evaluation configurations; only the \texttt{B-*} rows share an
implementation and support controlled attribution.}
\label{tab:configs}
\end{table*}

\wgpu{} is reported as an end-to-end tracked baseline because its validation,
lifetime, binding, and command models differ from \blade. CPU profiles identify
time spent in its trackers, while Vulkan captures report the barriers and
layouts those trackers generate. Direct \texttt{wgpu-hal} measurements are
optional diagnostics for separating frontend work; they are not required for
the main comparison and are not presented as another tracking implementation.

\subsection{Workloads}

Synthetic workloads vary pass count, resource count, resource size, operation
intensity, and dependency density. Graph families include serial chains,
independent passes, bounded-width pipelines, fan-out/fan-in, and randomized
directed acyclic graphs. Operations cover compute-to-compute,
transfer-to-compute, compute-to-graphics, attachment-to-sampling, and
graphics-to-graphics dependencies.

The application set will include a draw-call stress test, a multi-pass
renderer, a particle simulation with compute-to-graphics dependencies, and a
compute pipeline representative of matrix or tensor processing.
\pending{Freeze exact revisions, scenes, resolutions, and input hashes.}

\subsection{Platforms}

The primary matrix includes at least two AMD and two NVIDIA architecture
generations where available. Each result records GPU and CPU model, operating
system, kernel, Vulkan loader, driver, API version, enabled extensions, clocks,
power policy, and the exact repository revision.

Metal is not pooled into the Vulkan layout comparison. Its API hides image
layouts and may perform framework-level hazard tracking, so Apple results
require a separate tracked-versus-untracked methodology.

\subsection{Metrics and statistics}

We separately measure command recording, submission, host waiting, and GPU
execution. Driver profiling supplies barrier stalls, cache operations,
compression state where exposed, and overlap on the GPU timeline. Memory use
and barrier counts are also recorded.

Each configuration is warmed before measurement and sampled repeatedly in a
counterbalanced order. We report medians, distributions, and bootstrap
confidence intervals rather than a subjective workload threshold. Validation
layers are enabled for correctness runs and disabled for performance runs, as
recommended by vendor guidance~\cite{nvidia-vulkan-tips}.

\section{Results}

\subsection{Host-side scaling}
\pending{Plot recording time against passes and resources; attribute tracker
merge and transition-generation cost separately.}

\subsection{Barrier placement and dependency density}
\pending{Plot automatic-pass versus hazard-only global barriers over graph
width and dependency density. Include GPU timelines for representative
crossovers.}

\subsection{Matched \blade{} versus \wgpu{}}
\pending{Report end-to-end host and GPU differences. Use tracker profiles and
captured command streams as explanatory evidence, without attributing the
whole difference to tracking alone.}

\subsection{Generated barriers and image layouts}
\pending{Compare captured barrier counts, scopes, and image layouts. Stratify
observed performance by unified-image-layout support, image use, format, MSAA,
and architecture generation; do not claim independent layout causality.}

\subsection{Application workloads}
\pending{Report CPU and GPU deltas with absolute frame or job times. Explain
which synthetic regime predicts each application result.}

\section{Discussion}

The final analysis should produce a decision boundary, not a universal winner.
Expected explanatory variables include dependency density, opportunities for
stage overlap, the number of resources touched per pass, image compression
behavior, and whether synchronization is CPU- or GPU-critical.

Three outcomes are all actionable. First, negligible device cost with
measurable host savings would support a tracking-free contract for constrained
workloads. Second, architecture-specific layout or cache penalties would
motivate capability-gated policies. Third, losses on wide graphs but not serial
chains would suggest placing a render-graph or scheduling layer above a coarse
low-level API rather than requiring the low-level API itself to track states.

\section{Threats to Validity}

GPU timestamps and profiling can perturb short workloads. We therefore compare
instrumented measurements with externally captured timelines and amortized
throughput runs. Driver heuristics may recognize repeated synthetic work, so
inputs and resources vary while preserving the graph shape. Results from one
driver version cannot establish permanent vendor behavior; all raw metadata is
retained.

The controlled \blade{} variants isolate barrier placement but not global
versus resource-specific scope or persistent \texttt{GENERAL} versus
usage-specific layouts. Those dimensions are observed through the matched
\wgpu{} comparison and captured command streams, so causal claims about them
are deliberately limited. End-to-end \wgpu{} differences cannot be attributed
solely to tracking and are labeled as such.

\section{Related Work}

The Vulkan specification defines explicit execution and memory dependencies
and application-controlled image layouts~\cite{vulkan-spec}. AMD and NVIDIA
have published architecture-specific synchronization guidance
~\cite{amd-barriers,amd-rdna-guide,nvidia-vulkan-tips}. The
unified-image-layouts extension standardizes the observation that specialized
layouts no longer correspond to distinct representations on many modern
implementations~\cite{vulkan-unified-layouts}.

WebGPU specifies portable usage-scope validation~\cite{webgpu-spec}; \wgpu{}
implements it using optimized dense resource trackers~\cite{wgpu-trackers}.
\blade{} occupies a different point in this design space by exposing an unsafe
native contract and relying on application-declared pass dependencies.

\section{Conclusion}

\pending{State the measured workload envelope, quantified trade-offs, and
recommendations. Do not conclude that tracking is universally unnecessary.}

\bibliographystyle{plain}
\bibliography{references}

\end{document}
